\documentclass{article}
\usepackage{neurips_2024}

\usepackage[utf8]{inputenc}
\usepackage[T1]{fontenc}
\usepackage{hyperref}
\usepackage{url}
\usepackage{booktabs}
\usepackage{amsfonts}
\usepackage{amsmath}
\usepackage{amssymb}
\usepackage{amsthm}
\usepackage{nicefrac}
\usepackage{microtype}
\usepackage{xcolor}
\usepackage{algorithm}
\usepackage{algorithmic}
\usepackage{enumitem}
\usepackage{tikz}
\usepackage{stmaryrd}

% Theorem environments
\newtheorem{definition}{Definition}
\newtheorem{assumption}{Assumption}
\newtheorem{proposition}{Proposition}
\newtheorem{theorem}{Theorem}
\newtheorem{lemma}{Lemma}
\newtheorem{corollary}{Corollary}
\newtheorem{remark}{Remark}
\newtheorem{example}{Example}

% Operator shortcuts
\DeclareMathOperator*{\argmin}{arg\,min}
\DeclareMathOperator*{\argmax}{arg\,max}
\DeclareMathOperator{\pre}{pre}
\DeclareMathOperator{\eff}{eff}
\DeclareMathOperator{\dom}{dom}
\DeclareMathOperator{\img}{img}
\newcommand{\calS}{\mathcal{S}}
\newcommand{\calA}{\mathcal{A}}
\newcommand{\calO}{\mathcal{O}}
\newcommand{\calC}{\mathcal{C}}
\newcommand{\calD}{\mathcal{D}}
\newcommand{\calE}{\mathcal{E}}
\newcommand{\calF}{\mathcal{F}}
\newcommand{\calG}{\mathcal{G}}
\newcommand{\calH}{\mathcal{H}}
\newcommand{\calI}{\mathcal{I}}
\newcommand{\calK}{\mathcal{K}}
\newcommand{\calL}{\mathcal{L}}
\newcommand{\calM}{\mathcal{M}}
\newcommand{\calN}{\mathcal{N}}
\newcommand{\calP}{\mathcal{P}}
\newcommand{\calR}{\mathcal{R}}
\newcommand{\calT}{\mathcal{T}}
\newcommand{\calU}{\mathcal{U}}
\newcommand{\calV}{\mathcal{V}}
\newcommand{\calW}{\mathcal{W}}
\newcommand{\calX}{\mathcal{X}}
\newcommand{\bbR}{\mathbb{R}}
\newcommand{\bbN}{\mathbb{N}}
\newcommand{\bbZ}{\mathbb{Z}}
\newcommand{\ind}{\mathbb{1}}
\newcommand{\Obj}{\mathsf{Obj}}
\newcommand{\Pred}{\mathsf{Pred}}
\newcommand{\strips}{\textsc{Strips}}
\newcommand{\pddl}{\textsc{Pddl}}
\newcommand{\vlm}{\textsc{Vlm}}

\title{Object-Oriented World Model Learning from Raw Observations:\\A Formal Framework}

\author{}

\begin{document}

\maketitle

\begin{abstract}
We formalize object-oriented world model learning for interactive environments
observed as raw integer grids.  The agent must simultaneously discover a state
representation (object types, predicates, roles) and learn a transition model
(how the world changes under actions), starting from nothing.  We structure
both the representation and the model as an OOP class hierarchy: objects are
instances of learned classes with polymorphic update methods, and actions are
classes with polymorphic precondition and effect methods.  The class hierarchy
serves as the hypothesis space.  We extend WorldCoder's sample complexity
bound to this joint setting and show that the OOP factorization yields a
tighter bound than monolithic program synthesis.
\end{abstract}


% ============================================================================
% SECTION 1: ENVIRONMENT
% ============================================================================

\section{Environment}
\label{sec:environment}

\begin{definition}[Interactive Grid Environment]
\label{def:env}
An \textbf{interactive grid environment} is a deterministic MDP without reward:
\begin{equation}
\calM = (\calS, \calA, T, s_0),
\end{equation}
where:
\begin{itemize}[nosep]
    \item $\calS \subseteq \bbZ_K^{H \times W}$ is the set of reachable states.
          Each state $s \in \calS$ is an integer grid with $K$ colors:
          $s_{ij} \in \{0, \ldots, K{-}1\}$.
    \item $\calA = \{a_1, \ldots, a_N\}$ is a finite action set.  The agent
          knows $|\calA|$ but not the semantics of any action.
    \item $T : \calS \times \calA \to \calS$ is the true (unknown) deterministic
          transition function.
    \item $s_0 \in \calS$ is the initial state.
\end{itemize}
The agent observes $s_0$ and interacts by selecting actions, receiving the
resulting state after each action.  There is no reward signal during
exploration; the objective is model acquisition.
\end{definition}

\begin{definition}[Replay Buffer]
\label{def:replay}
After $t$ interactions, the agent's \textbf{replay buffer} is:
\begin{equation}
\calD_t = \{(s_0, a_0, s_1), (s_1, a_1, s_2), \ldots, (s_{t-1}, a_{t-1}, s_t)\}.
\end{equation}
\end{definition}

\begin{definition}[Environment Diameter]
\label{def:diameter}
The \textbf{diameter} of $\calM$ is the maximum over all pairs of reachable
states of the shortest action sequence connecting them:
\begin{equation}
D_{\calS,\calA,T} = \max_{s, s' \in \calS_{\text{reach}}}
    \min_{p \text{ a path from } s \text{ to } s'} \text{len}(p).
\end{equation}
\end{definition}


% ============================================================================
% SECTION 2: OOP STATE REPRESENTATION
% ============================================================================

\section{Object-Oriented State Representation}
\label{sec:oop_state}

The key departure from WorldCoder is that we structure both state and
transitions using object-oriented programming.  The world is composed of
typed objects.  Each object type is a class in a hierarchy.  Transition
logic is distributed across objects via polymorphic methods.

\subsection{Object Detection}
\label{sec:detection}

\begin{definition}[Raw Object]
\label{def:raw_object}
A \textbf{raw object} extracted from grid state $s$ is a tuple $o = (c, P)$,
where $c \in \{1, \ldots, K{-}1\}$ is the color (excluding background color 0)
and $P \subseteq [H] \times [W]$ is a connected set of cells under
8-connectivity.
\end{definition}

\begin{definition}[Object Detection Function]
\label{def:detection}
The \textbf{object detection function} $\Omega : \calS \to 2^{\calO_{\text{raw}}}$
extracts raw objects via connected-component labeling on each non-background
color channel:
\begin{equation}
\Omega(s) = \bigsqcup_{c=1}^{K-1} \textsc{CC}_8(\{(i,j) : s_{ij} = c\}).
\end{equation}
This function is deterministic, parameter-free, and domain-agnostic.
\end{definition}

\subsection{Object Type Hierarchy}
\label{sec:type_hierarchy}

Each raw object is an instance of some \emph{object type} (class) in a
hierarchy rooted at a generic base class.

\begin{definition}[Object Type]
\label{def:object_type}
An \textbf{object type} (class) $\tau$ is a tuple:
\begin{equation}
\tau = (\text{name}, \; \tau_{\text{parent}}, \; \calP_\tau^{\text{atomic}},
       \; \calP_\tau^{\text{derived}}, \; f_\tau^{\text{respond}}),
\end{equation}
where:
\begin{itemize}[nosep]
    \item $\text{name} \in \Sigma^*$ is the class name
          (e.g., \texttt{Wall}, \texttt{Player}, \texttt{PushableBlock}).
    \item $\tau_{\text{parent}}$ is the parent type (or $\bot$ for the root
          class \texttt{GameObject}).
    \item $\calP_\tau^{\text{atomic}}$ is the set of \textbf{atomic predicates}
          (fluent properties stored explicitly as mutable state):
          e.g., $\mathsf{position}$, $\mathsf{color}$, $\mathsf{alive}$.
    \item $\calP_\tau^{\text{derived}}$ is the set of \textbf{derived predicates}
          (computed from atomic predicates via Python functions, never stored):
          e.g., $\mathsf{can\_move}()$, $\mathsf{is\_blocked}()$,
          $\mathsf{is\_adjacent}(o')$.
    \item $f_\tau^{\text{respond}} : (\text{self}, \text{action}, \text{state})
          \to \text{self}'$ is the \textbf{response method}: given an action and
          the current world state, compute this object's updated state.  This
          is the polymorphic method that subclasses override.
\end{itemize}
\end{definition}

\begin{definition}[Type Hierarchy]
\label{def:type_tree}
The \textbf{type hierarchy} $\calH_\tau = (\calV, E)$ is a rooted tree where:
\begin{itemize}[nosep]
    \item $\calV$ is the set of all object types (classes).
    \item $E = \{(\tau, \tau_{\text{parent}}) : \tau \in \calV,\;
          \tau_{\text{parent}} \neq \bot\}$ are directed edges from child to
          parent.
    \item The root is \texttt{GameObject}, the abstract base class.
    \item Inheritance: a child type $\tau_c$ inherits all atomic predicates,
          derived predicates, and methods from its parent $\tau_p$, and may
          override any of them.
\end{itemize}
\end{definition}

\begin{example}[Sokoban-like hierarchy]
\label{ex:hierarchy}
\begin{verbatim}
GameObject (root)
  +-- StaticObject
  |     +-- Wall
  |     +-- Goal
  +-- DynamicObject
        +-- Player
        +-- PushableBlock
              +-- HeavyBlock (override: requires two pushes)
\end{verbatim}
\texttt{StaticObject} overrides $f^{\text{respond}}$ to be a no-op (static
objects do not change).  \texttt{PushableBlock} overrides
$f^{\text{respond}}$ to move in the push direction if unblocked.
\texttt{HeavyBlock} further overrides to require a second push.
\end{example}

\subsection{Atomic vs.\ Derived Predicates}
\label{sec:atomic_derived}

\begin{definition}[Atomic Predicate]
\label{def:atomic}
An \textbf{atomic predicate} is a property of an object (or pair of objects)
that is stored explicitly as part of the mutable state.  Only atomic predicates
are modified by action effects.  Examples: $\mathsf{position}(o)$,
$\mathsf{color}(o)$, $\mathsf{alive}(o)$, $\mathsf{connected}(o_1, o_2)$.
\end{definition}

\begin{definition}[Derived Predicate]
\label{def:derived}
A \textbf{derived predicate} is a property computed from atomic predicates
via a deterministic function.  Derived predicates are never stored; they are
recomputed on demand.  Examples:
\begin{align}
\mathsf{can\_move}(o) &= \neg \exists o' \in \text{state} :
    \mathsf{is\_blocking}(o', o), \\
\mathsf{can\_reach}(o_1, o_2) &= \text{transitive closure of }
    \mathsf{connected}, \\
\mathsf{is\_isolated}(o) &= \neg \exists o' \neq o :
    \mathsf{adjacent}(o, o').
\end{align}
In Python, derived predicates are implemented as methods or properties on the
object class.
\end{definition}

\begin{remark}[Why the distinction matters]
\label{rem:atomic_derived}
Only atomic predicates appear in the state representation.  This keeps the
state minimal and avoids redundancy: after an action updates atomic predicates,
all derived predicates are automatically consistent (they are recomputed from
the updated atomic state).  This eliminates the ``frame problem'' for derived
predicates.  It is analogous to the distinction between fluents and derived
predicates in PDDL, and between base and view tables in databases.
\end{remark}

\subsection{Typed Object Instance and Structured State}
\label{sec:structured_state}

\begin{definition}[Typed Object Instance]
\label{def:typed_instance}
A \textbf{typed object instance} is a pair $(o, \tau)$ where $o$ is a raw
object and $\tau \in \calV$ is its assigned type.  The instance carries:
\begin{itemize}[nosep]
    \item A unique identifier (track ID) $\text{id}(o)$ for cross-timestep
          tracking.
    \item Atomic predicate values: the current value of each predicate in
          $\calP_\tau^{\text{atomic}}$.
    \item Access to all derived predicates and the response method via
          the class $\tau$.
\end{itemize}
\end{definition}

\begin{definition}[Structured State]
\label{def:structured_state}
The \textbf{structured state} at time $t$ is:
\begin{equation}
\bar{s}_t = (s_t, \; \{(o_i, \tau_i, \mathbf{p}_i)\}_{i=1}^{M_t}),
\end{equation}
where $s_t$ is the raw grid, $\{o_i\}$ are the detected objects,
$\{\tau_i\}$ are their assigned types, and $\{\mathbf{p}_i\}$ are
their atomic predicate values.  The structured state is computed from the
raw state by the \textbf{perception function}:
\begin{equation}
\text{Perceive}(s_t) = \text{TypeAssign}(\text{Track}(\Omega(s_t), \bar{s}_{t-1})),
\end{equation}
where $\text{Track}$ matches objects across frames (by color, position overlap,
and nearest-neighbor) and $\text{TypeAssign}$ assigns types from the current
hierarchy $\calH_\tau$.
\end{definition}

\subsection{Binary Relations and Ownership}
\label{sec:relations}

\begin{definition}[Relation Ownership]
\label{def:relation_ownership}
Each atomic binary relation $R(o_1, o_2)$ is \textbf{owned by} both
participating objects.  For symmetric relations (e.g., $\mathsf{connected}$),
the relation is stored once and both objects have read access.  For asymmetric
relations (e.g., $\mathsf{pushing}(o_1, o_2)$), the source object owns the
relation.  When object $o_1$ updates in response to an action, it may create,
modify, or remove relations it owns.
\end{definition}


% ============================================================================
% SECTION 3: OOP ACTION REPRESENTATION
% ============================================================================

\section{Object-Oriented Action Representation}
\label{sec:oop_actions}

Actions are also organized as a class hierarchy.

\begin{definition}[Action Type]
\label{def:action_type}
An \textbf{action type} (class) $\alpha$ is a tuple:
\begin{equation}
\alpha = (\text{name}, \; \alpha_{\text{parent}}, \;
          f_\alpha^{\text{pre}}, \; f_\alpha^{\text{apply}}),
\end{equation}
where:
\begin{itemize}[nosep]
    \item $\text{name} \in \Sigma^*$ is the action class name
          (e.g., \texttt{MoveUp}, \texttt{Push}, \texttt{Rotate}).
    \item $\alpha_{\text{parent}}$ is the parent action type (or $\bot$ for
          the root class \texttt{Action}).
    \item $f_\alpha^{\text{pre}} : \bar{s} \to \{0, 1\}$ is the
          \textbf{precondition method}: returns whether this action is
          applicable in structured state $\bar{s}$.
    \item $f_\alpha^{\text{apply}} : \bar{s} \to \bar{s}'$ is the
          \textbf{apply method}: computes the next structured state by
          invoking $f_\tau^{\text{respond}}$ on each affected object.
\end{itemize}
\end{definition}

\begin{definition}[Object-Centric Transition]
\label{def:oop_transition}
The \textbf{object-centric transition} for action $\alpha$ in structured state
$\bar{s} = (s, \{(o_i, \tau_i, \mathbf{p}_i)\})$ proceeds as follows:
\begin{enumerate}[nosep]
    \item \textbf{Precondition check}: if $f_\alpha^{\text{pre}}(\bar{s}) = 0$,
          return $\bar{s}$ (no-op).
    \item \textbf{Object response}: for each object $o_i$ in causal order,
          compute:
          \begin{equation}
          \mathbf{p}'_i = f_{\tau_i}^{\text{respond}}(o_i, \alpha, \bar{s}).
          \end{equation}
    \item \textbf{Physics}: apply action-independent effects (gravity,
          cascading) to all objects.
    \item \textbf{Render}: reconstruct the raw grid $s'$ from the updated
          objects.
\end{enumerate}
The ordering in step 2 follows causal dependency: the action's direct target
updates first, then objects affected by that update, and so on.  For
deterministic sequential domains (including ARC-AGI-3), this ordering is
well-defined.
\end{definition}

\begin{remark}[Connection to STRIPS]
\label{rem:strips_connection}
This object-centric transition subsumes STRIPS.  The precondition method
$f_\alpha^{\text{pre}}$ encodes $\pre(\alpha)$.  Each object's response method
$f_{\tau_i}^{\text{respond}}$ encodes the conditional effects of $\alpha$ on
object $o_i$.  The add and delete lists are implicit in the difference between
$\mathbf{p}_i$ and $\mathbf{p}'_i$.  The advantage over STRIPS is that
conditional effects are factored by object type (via polymorphism) rather
than enumerated in a flat list.
\end{remark}

\begin{remark}[Why not a monolithic transition function]
\label{rem:not_monolithic}
WorldCoder synthesizes a single function
$\hat{T}(s, a) \to s'$ with an if-elif chain dispatching on $a$.  This has
three problems: (1) it does not factor by object, so fixing one object's
behavior risks breaking others; (2) it does not support transfer, since
the entire function is domain-specific; (3) it grows linearly in the number
of object types and action types.  The OOP factorization addresses all three:
each class is an independent unit of synthesis, testing, and transfer.
\end{remark}


% ============================================================================
% SECTION 4: HYPOTHESIS SPACE
% ============================================================================

\section{The OOP Hierarchy as Hypothesis Space}
\label{sec:hypothesis_space}

The class hierarchy $\calH_\tau$ defines a structured hypothesis space for the
world model.  This is where we make our main technical contribution.

\subsection{World Model as Class Library}
\label{sec:class_library}

\begin{definition}[OOP World Model]
\label{def:oop_world_model}
An \textbf{OOP world model} is a tuple:
\begin{equation}
\hat{W} = (\calH_\tau, \; \calH_\alpha, \; f_{\text{perceive}}, \;
           f_{\text{render}}),
\end{equation}
where:
\begin{itemize}[nosep]
    \item $\calH_\tau$ is the object type hierarchy (Definition~\ref{def:type_tree}).
    \item $\calH_\alpha$ is the action type hierarchy.
    \item $f_{\text{perceive}} : \calS \to \bar{\calS}$ is the perception
          function (Definition~\ref{def:structured_state}).
    \item $f_{\text{render}} : \bar{\calS} \to \calS$ reconstructs a raw grid
          from a structured state.
\end{itemize}
The world model induces a transition function:
\begin{equation}
\hat{T}_{\hat{W}}(s, a) = f_{\text{render}}\bigl(
    f_a^{\text{apply}}(f_{\text{perceive}}(s))\bigr).
\end{equation}
\end{definition}

\subsection{Factored Model Class}
\label{sec:factored}

\begin{definition}[Factored Model Class]
\label{def:factored_class}
Let $\calV = \{\tau_1, \ldots, \tau_n\}$ be the object types and
$\calA_\alpha = \{\alpha_1, \ldots, \alpha_m\}$ be the action types.  The
\textbf{factored model class} is:
\begin{equation}
\calW = \calW_{\text{perceive}} \times
        \prod_{i=1}^{n} \calW_{\tau_i} \times
        \prod_{j=1}^{m} \calW_{\alpha_j} \times
        \calW_{\text{render}},
\end{equation}
where $\calW_{\tau_i}$ is the set of possible implementations of object type
$\tau_i$ (its response method and derived predicates) and $\calW_{\alpha_j}$
is the set of possible implementations of action type $\alpha_j$ (its
precondition and apply methods).
\end{definition}

\begin{proposition}[Factored Logical Dimensionality]
\label{prop:factored_dim}
The logical dimensionality of the factored model class satisfies:
\begin{equation}
K_{\bot\!\!\bot}(\calW) \leq \sum_{i=1}^{n} K_{\bot\!\!\bot}(\calW_{\tau_i})
    + \sum_{j=1}^{m} K_{\bot\!\!\bot}(\calW_{\alpha_j})
    + K_{\bot\!\!\bot}(\calW_{\text{perceive}}).
\end{equation}
That is, the total number of counterexamples needed is at most the sum of
counterexamples needed per component, rather than the product (which is what
a monolithic model class would give in the worst case).
\end{proposition}

\begin{proof}[Proof sketch]
A counterexample $(s, a, s')$ where $\hat{T}(s,a) \neq s'$ must be
attributable to an error in at least one component: either a specific object
type's response method produced the wrong output, or a specific action type's
precondition or apply method was wrong, or the perception function
misidentified an object.  Since errors are localized to specific components,
each counterexample eliminates at least one candidate implementation from one
component's hypothesis space.  The total number of such eliminations is bounded
by the sum of per-component logical dimensionalities.
\end{proof}

\subsection{Monotonic Pruning}
\label{sec:pruning}

The tree structure of the type hierarchy enables efficient pruning.

\begin{proposition}[Monotonic Pruning]
\label{prop:pruning}
If an object type $\tau$'s behavior (response method) is inconsistent with a
set of observations $\calD$, then all subtypes $\tau'$ that inherit $\tau$'s
response method without override are also inconsistent with $\calD$.
\end{proposition}

\begin{proof}
If $\tau'$ inherits $f_\tau^{\text{respond}}$ without override, then
$f_{\tau'}^{\text{respond}} = f_\tau^{\text{respond}}$, and so any observation
that falsifies $f_\tau^{\text{respond}}$ also falsifies
$f_{\tau'}^{\text{respond}}$.
\end{proof}

This means that when a parent class's behavior is shown wrong, we can prune
the entire subtree of children that inherit that behavior.  Children that
\emph{override} the relevant method are not pruned and must be tested
independently.

\subsection{Transfer via Inheritance}
\label{sec:transfer}

\begin{definition}[Minimal Subclass Specialization]
\label{def:minimal_spec}
When moving to a new domain, the agent attempts to reuse existing object types.
If type $\tau$ from a previous domain does not match the observations in the new
domain, the agent derives a subclass $\tau'$ with $\tau_{\text{parent}}' = \tau$
that overrides the minimal set of methods needed to explain the new observations.
\end{definition}

\begin{example}[Movable walls]
\label{ex:movable_wall}
In domain A, \texttt{Wall} has $f^{\text{respond}}$ that returns self unchanged
(walls are static).  In domain B, some walls can be pushed.  The agent derives
\texttt{MovableWall} from \texttt{Wall}, overriding only
$f^{\text{respond}}$ to handle push actions while inheriting all other
properties (color, blocking behavior, rendering).
\end{example}

\begin{proposition}[Transfer Sample Complexity]
\label{prop:transfer}
Let $\hat{W}$ be a correct OOP world model for domain A, and let
$\hat{W}'$ be the correct model for domain B.  If $\hat{W}'$ can be obtained
from $\hat{W}$ by adding $k$ new subclasses and overriding $r$ methods total,
then the sample complexity for learning $\hat{W}'$ given $\hat{W}$ is:
\begin{equation}
D_{\calS',\calA',T'} \times \left(\sum_{i=1}^{k}
    K_{\bot\!\!\bot}(\calW_{\tau'_i}) + r + 1\right),
\end{equation}
compared to $D_{\calS',\calA',T'} \times (K_{\bot\!\!\bot}(\calW') + 1)$ for
learning from scratch.
\end{proposition}


% ============================================================================
% SECTION 5: JOINT LEARNING PROBLEM
% ============================================================================

\section{Joint Learning: Representation and Model}
\label{sec:joint}

The agent must simultaneously learn the type hierarchy (what object types
exist, what properties they have) and the transition logic (how each type
responds to each action).  This extends the problem formulated in our
previous work.

\subsection{The Predicate Grammar}
\label{sec:grammar}

\begin{definition}[Predicate Grammar]
\label{def:grammar}
The \textbf{predicate grammar} $\Psi_t$ at time $t$ is the current set of
candidate predicates (both atomic and derived) the agent could use.  $\Psi_t$
evolves over time as the agent encounters new phenomena.

The grammar contains:
\begin{itemize}[nosep]
    \item Atomic predicates: $\Psi_t^{\text{atomic}} =
          \bigcup_{\tau \in \calV_t} \calP_\tau^{\text{atomic}}$.
    \item Derived predicates: $\Psi_t^{\text{derived}} =
          \bigcup_{\tau \in \calV_t} \calP_\tau^{\text{derived}}$.
\end{itemize}
Growing $\Psi_t$ corresponds to adding new object types or adding new
predicates to existing types.
\end{definition}

\begin{definition}[Adequate Type Hierarchy]
\label{def:adequate}
A type hierarchy $\calH_\tau$ is \textbf{adequate} for environment $\calM$
if there exists an assignment of types to objects and an implementation of
all response methods such that the induced transition function $\hat{T}$
correctly predicts all transitions in $\calM$.

Formally: $\calH_\tau$ is adequate iff $T^* \in \calW(\calH_\tau)$, where
$\calW(\calH_\tau)$ is the factored model class induced by $\calH_\tau$.
\end{definition}

\subsection{Two Types of Counterexamples}
\label{sec:counterexamples}

\begin{definition}[Counterexample Types]
\label{def:cx_types}
A counterexample is a transition $(s, a, s')$ where $\hat{T}(s, a) \neq s'$.
We distinguish:

\begin{enumerate}[nosep]
    \item \textbf{Model counterexample}: the type hierarchy $\calH_\tau$ is
          adequate, but the current method implementations are wrong.  There
          exists a better implementation within $\calW(\calH_\tau)$ that
          resolves the counterexample.  Fix: re-synthesize the relevant
          response or apply method.

    \item \textbf{Representation counterexample}: the type hierarchy
          $\calH_\tau$ is inadequate.  No implementation within
          $\calW(\calH_\tau)$ can explain the observation.  Fix: grow the
          hierarchy by adding new types, new predicates, or new subclass
          specializations.
\end{enumerate}
\end{definition}

\begin{proposition}[Distinguishing Counterexample Types]
\label{prop:distinguish}
Model counterexamples are resolvable by re-synthesis: after $k$ synthesis
attempts (for bounded $k$), a correct implementation is found.
Representation counterexamples persist across synthesis attempts: no
implementation within the current model class resolves them.

Operationally: if synthesis fails to resolve a counterexample after $k_{\max}$
attempts, declare it a representation counterexample and trigger hierarchy
refinement.
\end{proposition}

\subsection{The Learning Algorithm}
\label{sec:algorithm}

The agent runs two nested loops.

\begin{algorithm}[h]
\caption{OOP World Model Learning}
\label{alg:main}
\begin{algorithmic}[1]
\STATE \textbf{Initialize:} $\calH_\tau \leftarrow \{\texttt{GameObject}\}$,
       $\calH_\alpha \leftarrow \{\texttt{Action}\}$, $\calD \leftarrow \emptyset$
\REPEAT
    \STATE \textbf{// Inner loop: WorldCoder-style exploration + synthesis}
    \REPEAT
        \STATE Select action $a_t$ via exploration policy (Section~\ref{sec:exploration})
        \STATE Observe transition $(s_t, a_t, s_{t+1})$;
               add to $\calD$
        \STATE Compute structured diff between
               $\text{Perceive}(s_t)$ and $\text{Perceive}(s_{t+1})$
        \STATE Check $\hat{T}(s_t, a_t)$ against $s_{t+1}$
        \IF{prediction incorrect}
            \STATE Diagnose: which object type or action type is wrong?
            \STATE Attempt re-synthesis of the faulty component
            \STATE Regression-test against $\calD$; accept if no regressions
        \ENDIF
        \STATE Track persistent counterexamples (not resolved after
               $k_{\max}$ synthesis attempts)
    \UNTIL{all counterexamples resolved \textbf{or}
           persistent counterexamples detected}
    \STATE \textbf{// Outer loop: hierarchy refinement}
    \IF{persistent (representation) counterexamples exist}
        \STATE Propose new object types or predicates based on the nature of
               the unresolved counterexamples
        \STATE Grow $\calH_\tau$ by adding subclasses or new branches
        \STATE Re-assign object types using the expanded hierarchy
    \ENDIF
\UNTIL{model is correct on all of $\calD$}
\end{algorithmic}
\end{algorithm}

\subsection{Sample Complexity}
\label{sec:sample_complexity}

\begin{theorem}[Joint Learning Bound]
\label{thm:joint}
Let $\calH_\tau^*$ be the true adequate type hierarchy and $\hat{W}^*$ be
the true world model.  Suppose the agent uses optimistic exploration
(WorldCoder's $\phi_2$) and the synthesis oracle succeeds within $k_{\max}$
attempts for model counterexamples.  Then the total number of environment
interactions to learn a correct world model is at most:
\begin{equation}
D_{\calS,\calA,T} \times \left(
    \underbrace{|\Psi_{\text{total}}|}_{\text{representation cost}}
    + \underbrace{\sum_{\tau \in \calH_\tau^*}
    K_{\bot\!\!\bot}(\calW_\tau)
    + \sum_{\alpha \in \calH_\alpha^*}
    K_{\bot\!\!\bot}(\calW_\alpha)}_{\text{model cost}}
    + 1
\right),
\end{equation}
where $|\Psi_{\text{total}}| = |\Psi_0 \cup \Psi_1 \cup \cdots \cup \Psi_T|$
is the total number of distinct predicates ever considered across all outer
loop iterations.
\end{theorem}

\begin{proof}[Proof sketch]
The outer loop iterates at most $|\Psi_{\text{total}}|$ times, since each
iteration either adds a new predicate/type (drawn from the finite cumulative
grammar) or the hierarchy is already adequate.  Each outer loop iteration
runs the inner loop, which by WorldCoder's Theorem~2.4 requires at most
$D \times (K_{\bot\!\!\bot}(\calW(\calH_\tau^{(i)})) + 1)$ interactions,
where $\calH_\tau^{(i)}$ is the hierarchy at outer iteration $i$.  By
Proposition~\ref{prop:factored_dim}, the factored logical dimensionality is
at most the sum of per-component dimensionalities.  Summing across outer
iterations and noting that each component's counterexamples are cumulative
(they carry over and don't need to be re-discovered) gives the bound.
\end{proof}

\begin{corollary}[Advantage of OOP Factorization]
\label{cor:advantage}
For a monolithic model class $\calW_{\text{mono}}$ with
$K_{\bot\!\!\bot}(\calW_{\text{mono}})$ logical dimensionality, the
OOP-factored bound replaces $K_{\bot\!\!\bot}(\calW_{\text{mono}})$ with
$\sum_i K_{\bot\!\!\bot}(\calW_{\tau_i}) + \sum_j K_{\bot\!\!\bot}(\calW_{\alpha_j})$.
Since each component's hypothesis space is strictly smaller than the monolithic
space, and logical dimensionality is subadditive under Cartesian product, the
factored bound is tighter.
\end{corollary}


% ============================================================================
% SECTION 6: PERCEPTION PIPELINE
% ============================================================================

\section{Perception Pipeline}
\label{sec:perception}

The perception pipeline bridges raw grid observations to the structured OOP
state.

\subsection{Object Detection and Tracking}
\label{sec:tracking}

Object detection uses connected-component labeling
(Definition~\ref{def:detection}).  Object \textbf{tracking} across frames
uses a three-pass matching algorithm:
\begin{enumerate}[nosep]
    \item \textbf{Color match}: objects with the same color.
    \item \textbf{Spatial overlap}: among same-color objects, match by
          maximum pixel overlap between frames.
    \item \textbf{Nearest neighbor}: unmatched objects are assigned to the
          nearest same-color object by centroid distance.
\end{enumerate}
Objects that appear in $s_{t+1}$ but not $s_t$ are new (appeared); objects
in $s_t$ but not $s_{t+1}$ have disappeared.

\subsection{Type Assignment}
\label{sec:type_assignment}

Given detected and tracked objects, the agent assigns types from the current
hierarchy.  The VLM proposes type assignments based on visual appearance,
behavior under actions, and context.  Type assignment is revisable: if an
object's behavior contradicts its assigned type, the assignment is updated.

\begin{definition}[Type Assignment Function]
\label{def:type_assign}
The \textbf{type assignment function} at time $t$ is:
\begin{equation}
\text{TypeAssign}_t : \calO_{\text{raw}} \times \calD_t \times \calH_\tau
\to \calV,
\end{equation}
mapping each raw object to a type in the current hierarchy, conditioned on
the replay buffer (which reveals how objects have behaved under actions).
\end{definition}


% ============================================================================
% SECTION 7: PROGRAM SYNTHESIS
% ============================================================================

\section{Program Synthesis via LLM}
\label{sec:synthesis}

The LLM synthesizes Python code for each component of the OOP world model.

\subsection{What the LLM Synthesizes}
\label{sec:what_synth}

The LLM is asked to produce:
\begin{enumerate}[nosep]
    \item \textbf{Object type classes}: for each identified object type,
          a Python class inheriting from \texttt{GameObject} (or a more
          specific parent), with:
          \begin{itemize}[nosep]
              \item Atomic predicate declarations (instance variables).
              \item Derived predicate implementations (methods/properties).
              \item The \texttt{respond\_to\_action} method.
          \end{itemize}
    \item \textbf{Action type classes}: for each identified action, a Python
          class inheriting from \texttt{Action}, with:
          \begin{itemize}[nosep]
              \item The \texttt{preconditions\_met} method.
              \item The \texttt{apply} method, which calls
                    \texttt{respond\_to\_action} on affected objects.
          \end{itemize}
    \item \textbf{Physics rules}: an optional \texttt{apply\_physics} function
          for action-independent effects (gravity, cascading).
    \item \textbf{Observation function}: a \texttt{perceive} function that
          extracts typed objects from a raw grid state.
\end{enumerate}

\subsection{Abstract Base Classes}
\label{sec:abc}

The agent provides domain-generic abstract base classes that define the
interface contract.  The LLM fills in domain-specific implementations.

\begin{verbatim}
class GameObject(ABC):
    """Abstract base class for all game objects."""
    id: str
    position: tuple[int, int]
    color: int
    pixels: set[tuple[int, int]]

    @abstractmethod
    def respond_to_action(self, action: 'Action',
                          state: 'WorldState') -> None:
        """Update self in response to the given action."""

    # Derived predicates (override in subclasses)
    def can_move(self) -> bool: return True
    def is_blocking(self, other: 'GameObject') -> bool: return False

class Action(ABC):
    """Abstract base class for all actions."""
    name: str

    @abstractmethod
    def preconditions_met(self, state: 'WorldState') -> bool: ...

    @abstractmethod
    def apply(self, state: 'WorldState') -> 'WorldState':
        """Apply action: calls respond_to_action on each affected
        object in causal order."""

class WorldState:
    """The structured world state."""
    objects: list[GameObject]
    grid: np.ndarray  # raw grid for rendering

    def get_objects_of_type(self, cls: type) -> list[GameObject]: ...
    def get_object_at(self, pos: tuple[int, int]) -> GameObject | None: ...
    def get_adjacent(self, obj: GameObject) -> list[GameObject]: ...
\end{verbatim}

\subsection{Counterexample-Guided Refinement}
\label{sec:cegis}

Synthesis follows a counterexample-guided inductive synthesis (CEGIS) loop:

\begin{enumerate}[nosep]
    \item \textbf{Synthesize}: LLM produces class implementations given
          the current counterexamples, the abstract base classes, and the
          current NL descriptions of object/action behavior.
    \item \textbf{Verify}: run the synthesized model against the entire
          replay buffer $\calD$.  Collect mismatches.
    \item \textbf{Refine}: if mismatches exist, present them to the LLM
          along with the faulty code.  The LLM produces a corrected version.
    \item \textbf{Regression gate}: accept the new version only if it does
          not regress on previously correct predictions.
\end{enumerate}

The error signal is structured: for each mismatch, the agent reports which
object(s) predicted incorrectly, what the predicted vs.\ actual state was,
and which method was responsible.  This per-object, per-method attribution
enables targeted fixes rather than wholesale rewrites.


% ============================================================================
% SECTION 8: EXPLORATION
% ============================================================================

\section{Active Exploration}
\label{sec:exploration}

\subsection{Exploration Policy}
\label{sec:explore_policy}

The exploration policy selects actions to maximize information gain about the
world model.  Following WorldCoder, the agent maintains an optimism constraint
($\phi_2$): the current model must be consistent with some path from the
current state to positive reward (or, in our reward-free setting, to unvisited
state regions).

The policy prioritizes:
\begin{enumerate}[nosep]
    \item \textbf{Untested action-object combinations}: actions that have not
          been observed to affect a particular object type.
    \item \textbf{Unverified object behaviors}: object types whose response
          methods have fewer than $\kappa$ consistent confirmations.
    \item \textbf{State diversity}: reaching states where objects are in
          novel configurations (different relative positions, different
          adjacency patterns).
    \item \textbf{Hypothesis testing}: targeted actions designed to
          distinguish between competing type assignments or method
          implementations.
\end{enumerate}

\subsection{Setup Sequences}
\label{sec:setup}

To reach informative states, the agent constructs \textbf{setup sequences}
of verified actions.  Only actions whose behavior has been confirmed by
$\kappa$ consecutive consistent observations are used in setup sequences,
providing a safety guarantee analogous to the safe action models of
Juba et al.

\subsection{Data Consistency and Optimism}
\label{sec:phi}

Following WorldCoder, the learned model $\hat{W}$ must satisfy:
\begin{align}
\phi_1(\calD, \hat{W}) &= \forall (s, a, s') \in \calD :
    \hat{T}_{\hat{W}}(s, a) = s'
    \tag{data consistency} \\
\phi_2(s_0, \hat{W}) &= \exists a_1, \ldots, a_\ell :
    \text{the trajectory under } \hat{T}_{\hat{W}} \text{ from } s_0
    \text{ reaches an unvisited region}
    \tag{optimism}
\end{align}
The agent plans over $\hat{T}_{\hat{W}}$ satisfying $\phi_1 \wedge \phi_2$,
using MCTS with the synthesized Python model as the simulator.


% ============================================================================
% SECTION 9: CONTINUAL LEARNING / TRANSFER
% ============================================================================

\section{Continual Learning and Cross-Domain Transfer}
\label{sec:continual}

The OOP structure naturally supports continual learning.  When moving to a
new domain:

\begin{enumerate}[nosep]
    \item \textbf{Reuse}: attempt to reuse the existing type hierarchy.
          Apply the current perception and transition model to the new domain.
    \item \textbf{Diagnose}: if counterexamples arise, determine which
          object types need modification.
    \item \textbf{Specialize}: derive new subclasses that override only
          the behavior that differs.
    \item \textbf{Generalize}: if a pattern recurs across domains, refactor
          common behavior into a shared parent class.
\end{enumerate}

This is a tree-structured version space: we start generic and specialize
as needed.  If an existing subclass matches, we reuse it.  If not, we
derive a new one.  Inconsistent subclasses (whose behavior contradicts
observations) are pruned along with their subtrees
(Proposition~\ref{prop:pruning}).

The concrete deliverable for the paper: design the abstract base classes
for \texttt{GameObject}, \texttt{Action}, \texttt{WorldState},
\texttt{ObservationFunction} in a domain-generic way.  Show they are
sufficient for a range of grid-world domains.  Show that new domains only
require subclass specializations.  Prove sample complexity bounds extending
WorldCoder's theorem to this structured OOP setting
(Theorem~\ref{thm:joint}).

\end{document}
